\question[5]
a. Schreibe eine Klasse \texttt{Node}, die einen Knoten eines binären Baums repräsentiert.

b. Schreibe eine Funktion \texttt{baum}, die den abgebildeten Baum zurückgibt.
Nutze dazu deine Node-Klasse.
Verschachtele die Node-Aufrufe aus Lesbarkeitsgründen
höchstens einmal und gib den temporären Teilbäumen sprechende Namen.

\includegraphics[height=4cm]{\pfad/baumdef_02.png}

\begin{solutionbox}{5.5cm}
\begin{lstlisting}
class Node:
    def __init__(self, val, left=None, right=None):
        self.val = val
        self.left = left
        self.right = right

def baum():
    h = Node('H',Node('B'),None)
    c = Node('C',h,None)
    a = Node('A',Node('D'),c)
    return a
\end{lstlisting}
\end{solutionbox}
